\begin{itemize}
    \item \textbf{CSU-01: Logar}
\end{itemize}

\textbf{Caso de uso:} Logar

\textbf{Identificador:} CSU-01

\textbf{Atores envolvidos:} Viewer, Editor

\textbf{Sumário:} Esta funcionalidade permite que o usuário previamente cadastrado autentique-se na plataforma, obtendo acesso às funcionalidades disponíveis.

\textbf{Pré-condição:}
\begin{itemize}
    \item O usuário deverá estar previamente cadastrado no sistema.
\end{itemize}

\textbf{Fluxo principal:}
\begin{itemize}
    \item O usuário acessa a tela de login;
    \item O usuário insere suas credenciais (e-mail e senha);
    \item O usuário submete as informações;
    \item O usuário é redirecionado para a tela de listagem de projetos (CSU-03).
\end{itemize}

\textbf{Exceção:}
\begin{itemize}
    \item \textbf{Credenciais inválidas:} Caso as credenciais informadas sejam inválidas, será mostrada uma mensagem de erro ao usuário.
    \item \textbf{Erro de instabilidade dos servidores:} Caso algum erro ocorra nos servidores, será apresentada uma mensagem explicando que houve alguma falha de comunicação no sistema e encorajando o usuário a tentar novamente mais tarde.
\end{itemize}

\textbf{Pós-Condições:}
\begin{itemize}
    \item O usuário poderá acessar as funcionalidades que o perfil de sua conta permitir.
\end{itemize}

\vspace{1em}

\begin{itemize}
    \item \textbf{CSU-03: Listar projetos}
\end{itemize}

\textbf{Caso de uso:} Listar projetos

\textbf{Identificador:} CSU-03

\textbf{Atores envolvidos:} Viewer, Editor

\textbf{Sumário:} Esta funcionalidade permite que o usuário autenticado visualize todos os projetos aos quais tem acesso como visualizador ou colaborador.

\textbf{Pré-condição:}
\begin{itemize}
    \item O usuário deverá estar autenticado no sistema (CSU-01).
\end{itemize}

\textbf{Fluxo principal:}
\begin{itemize}
    \item Após a autenticação (CSU-01), o usuário é redirecionado à tela de listagem de projetos;
    \item Os projetos aos quais o usuário tem acesso são exibidos;
    \item Caso deseje, o usuário poderá filtrar os projetos pelos campos disponíveis na página.
\end{itemize}

\textbf{Exceção:}
\begin{itemize}
    \item \textbf{Erro de instabilidade dos servidores:} Caso algum erro ocorra nos servidores, será apresentada uma mensagem explicando que houve alguma falha de comunicação no sistema e encorajando o usuário a tentar novamente mais tarde.
\end{itemize}

\textbf{Pós-Condições:}
\begin{itemize}
    \item O usuário poderá selecionar um projeto para visualizar seu diagrama ER (CSU-06).
\end{itemize}

\vspace{1em}

\begin{itemize}
    \item \textbf{CSU-06: Visualizar diagrama}
\end{itemize}

\textbf{Caso de uso:} Visualizar diagrama

\textbf{Identificador:} CSU-06

\textbf{Atores envolvidos:} Viewer, Editor

\textbf{Sumário:} Esta funcionalidade permite que o usuário visualize o diagrama ER de um projeto ao qual tem acesso, podendo observar as entidades, atributos e relacionamentos modelados.

\textbf{Pré-condição:}
\begin{itemize}
    \item O usuário deverá estar autenticado no sistema;
    \item O usuário deverá ter acesso ao projeto selecionado.
\end{itemize}

\textbf{Fluxo principal:}
\begin{itemize}
    \item O usuário acessa a tela de listagem de projetos (CSU-03);
    \item O usuário identifica o projeto cujo diagrama deseja visualizar;
    \item O usuário seleciona o projeto;
    \item O usuário é redirecionado à tela do diagrama ER do projeto;
    \item O diagrama é carregado e exibido ao usuário com todas as entidades, atributos e relacionamentos.
\end{itemize}

\textbf{Exceção:}
\begin{itemize}
    \item \textbf{Projeto não encontrado:} Caso o usuário tente acessar um projeto que não existe na plataforma, será apresentada uma mensagem de erro indicando o ocorrido.
    \item \textbf{Usuário sem acesso ao projeto:} Caso o usuário tente acessar um projeto ao qual não tem permissão, será apresentada uma mensagem de erro e o usuário será redirecionado à tela de listagem de projetos (CSU-03).
    \item \textbf{Erro de instabilidade dos servidores:} Caso algum erro ocorra nos servidores, será apresentada uma mensagem explicando que houve alguma falha de comunicação no sistema e encorajando o usuário a tentar novamente mais tarde.
\end{itemize}

\textbf{Pós-Condições:}
\begin{itemize}
    \item O usuário poderá visualizar o diagrama ER com todas as entidades, atributos e relacionamentos cadastrados para o projeto.
    \item Caso o usuário possua o perfil de Editor, poderá também editar o diagrama (CSU-07) e exportar o script DDL correspondente (CSU-23).
\end{itemize}

\vspace{1em}

\begin{itemize}
    \item \textbf{CSU-07: Editar diagrama}
\end{itemize}

\textbf{Caso de uso:} Editar diagrama

\textbf{Identificador:} CSU-07

\textbf{Atores envolvidos:} Editor

\textbf{Sumário:} Esta funcionalidade permite que o Editor realize alterações no diagrama ER de um projeto, podendo gerenciar tabelas, atributos e relacionamentos.

\textbf{Pré-condição:}
\begin{itemize}
    \item O usuário deverá estar autenticado no sistema;
    \item O usuário deverá ter o perfil de Editor no projeto selecionado;
    \item O diagrama do projeto deverá estar sendo visualizado (CSU-06).
\end{itemize}

\textbf{Fluxo principal:}
\begin{itemize}
    \item Ao acessar o diagrama (CSU-06), o Editor dispõe das opções de edição habilitadas na interface;
    \item O Editor realiza as operações desejadas, podendo adicionar tabelas (CSU-13), editar tabelas (CSU-14), remover tabelas (CSU-15), adicionar relacionamentos (CSU-16), remover relacionamentos (CSU-17), mover elementos (CSU-18) ou importar um script SQL (CSU-24);
    \item As alterações são persistidas em tempo real e propagadas aos demais usuários conectados ao projeto via \textit{WebSocket}.
\end{itemize}

\textbf{Exceção:}
\begin{itemize}
    \item \textbf{Usuário sem permissão de edição:} Caso o usuário não possua o perfil de Editor no projeto, as opções de edição não estarão disponíveis e será exibida uma mensagem de erro caso a ação seja tentada.
    \item \textbf{Erro de instabilidade dos servidores:} Caso algum erro ocorra nos servidores, será apresentada uma mensagem explicando que houve alguma falha de comunicação no sistema e encorajando o usuário a tentar novamente mais tarde.
\end{itemize}

\textbf{Pós-Condições:}
\begin{itemize}
    \item O diagrama ER do projeto estará atualizado com as alterações realizadas pelo Editor.
\end{itemize}

\vspace{1em}

\begin{itemize}
    \item \textbf{CSU-13: Adicionar tabela}
\end{itemize}

\textbf{Caso de uso:} Adicionar tabela

\textbf{Identificador:} CSU-13

\textbf{Atores envolvidos:} Editor

\textbf{Sumário:} Esta funcionalidade permite que o Editor adicione uma nova tabela (entidade) ao diagrama ER do projeto.

\textbf{Pré-condição:}
\begin{itemize}
    \item O usuário deverá estar em modo de edição do diagrama (CSU-07).
\end{itemize}

\textbf{Fluxo principal:}
\begin{itemize}
    \item O Editor seleciona a opção de adicionar uma nova tabela na interface de edição;
    \item Uma nova tabela é criada no canvas do diagrama com um nome padrão;
    \item O Editor define o nome da tabela;
    \item O Editor pode adicionar atributos à tabela (CSU-19) e visualizar a listagem dos atributos cadastrados (CSU-20);
    \item As alterações são persistidas.
\end{itemize}

\textbf{Exceção:}
\begin{itemize}
    \item \textbf{Nome de tabela inválido:} Caso o nome informado seja inválido ou esteja em branco, será exibida uma mensagem de erro ao usuário.
    \item \textbf{Erro de instabilidade dos servidores:} Caso algum erro ocorra nos servidores, será apresentada uma mensagem explicando que houve alguma falha de comunicação no sistema e encorajando o usuário a tentar novamente mais tarde.
\end{itemize}

\textbf{Pós-Condições:}
\begin{itemize}
    \item A nova tabela estará visível no diagrama ER do projeto.
\end{itemize}

\vspace{1em}

\begin{itemize}
    \item \textbf{CSU-14: Editar tabela}
\end{itemize}

\textbf{Caso de uso:} Editar tabela

\textbf{Identificador:} CSU-14

\textbf{Atores envolvidos:} Editor

\textbf{Sumário:} Esta funcionalidade permite que o Editor edite os dados de uma tabela existente no diagrama ER do projeto.

\textbf{Pré-condição:}
\begin{itemize}
    \item O usuário deverá estar em modo de edição do diagrama (CSU-07);
    \item A tabela a ser editada deverá existir no diagrama.
\end{itemize}

\textbf{Fluxo principal:}
\begin{itemize}
    \item O Editor seleciona uma tabela existente no diagrama;
    \item A interface de edição da tabela é exibida com os dados atuais;
    \item O Editor altera o nome ou outras propriedades da tabela;
    \item O Editor pode adicionar novos atributos (CSU-19) e gerenciar os atributos existentes por meio da listagem (CSU-20);
    \item As alterações são persistidas.
\end{itemize}

\textbf{Exceção:}
\begin{itemize}
    \item \textbf{Dados inválidos:} Caso algum dado informado seja inválido, será exibida uma mensagem de erro ao usuário.
    \item \textbf{Erro de instabilidade dos servidores:} Caso algum erro ocorra nos servidores, será apresentada uma mensagem explicando que houve alguma falha de comunicação no sistema e encorajando o usuário a tentar novamente mais tarde.
\end{itemize}

\textbf{Pós-Condições:}
\begin{itemize}
    \item A tabela estará atualizada no diagrama ER do projeto.
\end{itemize}

\vspace{1em}

\begin{itemize}
    \item \textbf{CSU-15: Remover tabela}
\end{itemize}

\textbf{Caso de uso:} Remover tabela

\textbf{Identificador:} CSU-15

\textbf{Atores envolvidos:} Editor

\textbf{Sumário:} Esta funcionalidade permite que o Editor remova uma tabela existente do diagrama ER do projeto.

\textbf{Pré-condição:}
\begin{itemize}
    \item O usuário deverá estar em modo de edição do diagrama (CSU-07);
    \item A tabela a ser removida deverá existir no diagrama.
\end{itemize}

\textbf{Fluxo principal:}
\begin{itemize}
    \item O Editor seleciona a tabela que deseja remover;
    \item O Editor aciona a opção de remoção;
    \item Uma janela de confirmação é exibida ao usuário;
    \item O Editor confirma a ação;
    \item A tabela e todos os relacionamentos a ela associados são removidos do diagrama;
    \item As alterações são persistidas.
\end{itemize}

\textbf{Exceção:}
\begin{itemize}
    \item \textbf{Tabela não encontrada:} Caso a tabela selecionada não exista no diagrama, será exibida uma mensagem de erro ao usuário.
    \item \textbf{Erro de instabilidade dos servidores:} Caso algum erro ocorra nos servidores, será apresentada uma mensagem explicando que houve alguma falha de comunicação no sistema e encorajando o usuário a tentar novamente mais tarde.
\end{itemize}

\textbf{Pós-Condições:}
\begin{itemize}
    \item A tabela e seus relacionamentos são removidos do diagrama ER do projeto.
\end{itemize}

\vspace{1em}

\begin{itemize}
    \item \textbf{CSU-16: Adicionar relacionamento}
\end{itemize}

\textbf{Caso de uso:} Adicionar relacionamento

\textbf{Identificador:} CSU-16

\textbf{Atores envolvidos:} Editor

\textbf{Sumário:} Esta funcionalidade permite que o Editor adicione um relacionamento entre duas tabelas no diagrama ER do projeto.

\textbf{Pré-condição:}
\begin{itemize}
    \item O usuário deverá estar em modo de edição do diagrama (CSU-07);
    \item Deverão existir pelo menos duas tabelas no diagrama.
\end{itemize}

\textbf{Fluxo principal:}
\begin{itemize}
    \item O Editor seleciona a opção de adicionar um relacionamento na interface de edição;
    \item O Editor seleciona a tabela de origem do relacionamento;
    \item O Editor seleciona a tabela de destino do relacionamento;
    \item O Editor define a cardinalidade do relacionamento;
    \item O relacionamento é criado e exibido no diagrama conectando as duas tabelas;
    \item As alterações são persistidas.
\end{itemize}

\textbf{Exceção:}
\begin{itemize}
    \item \textbf{Dados inválidos:} Caso algum dado informado seja inválido, será exibida uma mensagem de erro ao usuário.
    \item \textbf{Erro de instabilidade dos servidores:} Caso algum erro ocorra nos servidores, será apresentada uma mensagem explicando que houve alguma falha de comunicação no sistema e encorajando o usuário a tentar novamente mais tarde.
\end{itemize}

\textbf{Pós-Condições:}
\begin{itemize}
    \item O relacionamento estará visível no diagrama ER conectando as duas tabelas selecionadas.
\end{itemize}

\vspace{1em}

\begin{itemize}
    \item \textbf{CSU-17: Remover relacionamento}
\end{itemize}

\textbf{Caso de uso:} Remover relacionamento

\textbf{Identificador:} CSU-17

\textbf{Atores envolvidos:} Editor

\textbf{Sumário:} Esta funcionalidade permite que o Editor remova um relacionamento existente entre tabelas no diagrama ER do projeto.

\textbf{Pré-condição:}
\begin{itemize}
    \item O usuário deverá estar em modo de edição do diagrama (CSU-07);
    \item O relacionamento a ser removido deverá existir no diagrama.
\end{itemize}

\textbf{Fluxo principal:}
\begin{itemize}
    \item O Editor seleciona o relacionamento que deseja remover;
    \item O Editor aciona a opção de remoção;
    \item Uma janela de confirmação é exibida ao usuário;
    \item O Editor confirma a ação;
    \item O relacionamento é removido do diagrama;
    \item As alterações são persistidas.
\end{itemize}

\textbf{Exceção:}
\begin{itemize}
    \item \textbf{Relacionamento não encontrado:} Caso o relacionamento selecionado não exista no diagrama, será exibida uma mensagem de erro ao usuário.
    \item \textbf{Erro de instabilidade dos servidores:} Caso algum erro ocorra nos servidores, será apresentada uma mensagem explicando que houve alguma falha de comunicação no sistema e encorajando o usuário a tentar novamente mais tarde.
\end{itemize}

\textbf{Pós-Condições:}
\begin{itemize}
    \item O relacionamento é removido do diagrama ER do projeto.
\end{itemize}

\vspace{1em}

\begin{itemize}
    \item \textbf{CSU-18: Mover elemento no diagrama}
\end{itemize}

\textbf{Caso de uso:} Mover elemento no diagrama

\textbf{Identificador:} CSU-18

\textbf{Atores envolvidos:} Editor

\textbf{Sumário:} Esta funcionalidade permite que o Editor reposicione tabelas e demais elementos no canvas do diagrama ER do projeto.

\textbf{Pré-condição:}
\begin{itemize}
    \item O usuário deverá estar em modo de edição do diagrama (CSU-07);
    \item O elemento a ser movido deverá existir no diagrama.
\end{itemize}

\textbf{Fluxo principal:}
\begin{itemize}
    \item O Editor seleciona o elemento que deseja reposicionar no canvas;
    \item O Editor arrasta o elemento até a posição desejada;
    \item O elemento é exibido na nova posição;
    \item A nova posição é persistida.
\end{itemize}

\textbf{Exceção:}
\begin{itemize}
    \item \textbf{Erro de instabilidade dos servidores:} Caso algum erro ocorra nos servidores, será apresentada uma mensagem explicando que houve alguma falha de comunicação no sistema e encorajando o usuário a tentar novamente mais tarde.
\end{itemize}

\textbf{Pós-Condições:}
\begin{itemize}
    \item O elemento é exibido na nova posição no diagrama ER do projeto.
\end{itemize}

\vspace{1em}

\begin{itemize}
    \item \textbf{CSU-19: Adicionar atributo}
\end{itemize}

\textbf{Caso de uso:} Adicionar atributo

\textbf{Identificador:} CSU-19

\textbf{Atores envolvidos:} Editor

\textbf{Sumário:} Esta funcionalidade permite que o Editor adicione um novo atributo a uma tabela no diagrama ER do projeto.

\textbf{Pré-condição:}
\begin{itemize}
    \item O usuário deverá estar em modo de edição do diagrama (CSU-07);
    \item A tabela de destino deverá existir no diagrama.
\end{itemize}

\textbf{Fluxo principal:}
\begin{itemize}
    \item Ao adicionar (CSU-13) ou editar (CSU-14) uma tabela, o Editor seleciona a opção de adicionar um atributo;
    \item O Editor preenche as informações do atributo, como nome, tipo de dado, chave primária, chave estrangeira, restrição \textit{NOT NULL}, unicidade, auto incremento e valor padrão;
    \item O Editor submete os dados do atributo;
    \item O atributo é adicionado à tabela no diagrama;
    \item As alterações são persistidas.
\end{itemize}

\textbf{Exceção:}
\begin{itemize}
    \item \textbf{Dados inválidos:} Caso algum dado informado seja inválido, será exibida uma mensagem de erro indicando qual campo está incorreto.
    \item \textbf{Erro de instabilidade dos servidores:} Caso algum erro ocorra nos servidores, será apresentada uma mensagem explicando que houve alguma falha de comunicação no sistema e encorajando o usuário a tentar novamente mais tarde.
\end{itemize}

\textbf{Pós-Condições:}
\begin{itemize}
    \item O novo atributo estará listado na tabela no diagrama ER do projeto.
\end{itemize}

\vspace{1em}

\begin{itemize}
    \item \textbf{CSU-20: Listar atributos}
\end{itemize}

\textbf{Caso de uso:} Listar atributos

\textbf{Identificador:} CSU-20

\textbf{Atores envolvidos:} Editor

\textbf{Sumário:} Esta funcionalidade permite que o Editor visualize todos os atributos de uma tabela no diagrama ER do projeto.

\textbf{Pré-condição:}
\begin{itemize}
    \item O usuário deverá estar em modo de edição do diagrama (CSU-07);
    \item A tabela consultada deverá existir no diagrama.
\end{itemize}

\textbf{Fluxo principal:}
\begin{itemize}
    \item Ao adicionar (CSU-13) ou editar (CSU-14) uma tabela, os atributos cadastrados para aquela tabela são exibidos ao Editor;
    \item O Editor pode visualizar todos os atributos com suas respectivas configurações;
    \item Caso deseje, o Editor pode editar um atributo (CSU-21) ou remover um atributo (CSU-22).
\end{itemize}

\textbf{Exceção:}
\begin{itemize}
    \item \textbf{Erro de instabilidade dos servidores:} Caso algum erro ocorra nos servidores, será apresentada uma mensagem explicando que houve alguma falha de comunicação no sistema e encorajando o usuário a tentar novamente mais tarde.
\end{itemize}

\textbf{Pós-Condições:}
\begin{itemize}
    \item O Editor poderá gerenciar os atributos da tabela, podendo editá-los ou removê-los conforme necessário.
\end{itemize}

\vspace{1em}

\begin{itemize}
    \item \textbf{CSU-21: Editar atributo}
\end{itemize}

\textbf{Caso de uso:} Editar atributo

\textbf{Identificador:} CSU-21

\textbf{Atores envolvidos:} Editor

\textbf{Sumário:} Esta funcionalidade permite que o Editor edite as propriedades de um atributo existente em uma tabela do diagrama ER do projeto.

\textbf{Pré-condição:}
\begin{itemize}
    \item O usuário deverá estar em modo de edição do diagrama (CSU-07);
    \item O atributo a ser editado deverá existir na tabela.
\end{itemize}

\textbf{Fluxo principal:}
\begin{itemize}
    \item Ao listar os atributos de uma tabela (CSU-20), o Editor seleciona a opção de editar um atributo;
    \item A interface de edição do atributo é exibida com os valores atuais;
    \item O Editor altera os campos desejados;
    \item O Editor salva as alterações;
    \item O atributo é atualizado na tabela no diagrama;
    \item As alterações são persistidas.
\end{itemize}

\textbf{Exceção:}
\begin{itemize}
    \item \textbf{Dados inválidos:} Caso algum dado informado seja inválido, será exibida uma mensagem de erro indicando qual campo está incorreto.
    \item \textbf{Erro de instabilidade dos servidores:} Caso algum erro ocorra nos servidores, será apresentada uma mensagem explicando que houve alguma falha de comunicação no sistema e encorajando o usuário a tentar novamente mais tarde.
\end{itemize}

\textbf{Pós-Condições:}
\begin{itemize}
    \item O atributo estará atualizado na tabela no diagrama ER do projeto.
\end{itemize}

\vspace{1em}

\begin{itemize}
    \item \textbf{CSU-22: Remover atributo}
\end{itemize}

\textbf{Caso de uso:} Remover atributo

\textbf{Identificador:} CSU-22

\textbf{Atores envolvidos:} Editor

\textbf{Sumário:} Esta funcionalidade permite que o Editor remova um atributo de uma tabela no diagrama ER do projeto.

\textbf{Pré-condição:}
\begin{itemize}
    \item O usuário deverá estar em modo de edição do diagrama (CSU-07);
    \item O atributo a ser removido deverá existir na tabela.
\end{itemize}

\textbf{Fluxo principal:}
\begin{itemize}
    \item Ao listar os atributos de uma tabela (CSU-20), o Editor seleciona a opção de remover um atributo;
    \item Uma janela de confirmação é exibida ao usuário;
    \item O Editor confirma a ação;
    \item O atributo é removido da tabela no diagrama;
    \item As alterações são persistidas.
\end{itemize}

\textbf{Exceção:}
\begin{itemize}
    \item \textbf{Atributo não encontrado:} Caso o atributo selecionado não exista na tabela, será exibida uma mensagem de erro ao usuário.
    \item \textbf{Erro de instabilidade dos servidores:} Caso algum erro ocorra nos servidores, será apresentada uma mensagem explicando que houve alguma falha de comunicação no sistema e encorajando o usuário a tentar novamente mais tarde.
\end{itemize}

\textbf{Pós-Condições:}
\begin{itemize}
    \item O atributo é removido da tabela no diagrama ER do projeto.
\end{itemize}

\vspace{1em}

\begin{itemize}
    \item \textbf{CSU-23: Exportar DDL}
\end{itemize}

\textbf{Caso de uso:} Exportar DDL

\textbf{Identificador:} CSU-23

\textbf{Atores envolvidos:} Editor

\textbf{Sumário:} Esta funcionalidade permite que o Editor exporte o diagrama ER do projeto na forma de um script DDL (\textit{Data Definition Language}) em SQL.

\textbf{Pré-condição:}
\begin{itemize}
    \item O usuário deverá estar visualizando um diagrama (CSU-06);
    \item O diagrama deverá possuir ao menos uma tabela cadastrada.
\end{itemize}

\textbf{Fluxo principal:}
\begin{itemize}
    \item Ao visualizar o diagrama (CSU-06), o Editor seleciona a opção de exportar o script DDL;
    \item O sistema gera o script DDL correspondente ao estado atual do diagrama ER;
    \item O arquivo gerado é disponibilizado para \textit{download};
    \item O Editor realiza o \textit{download} do arquivo.
\end{itemize}

\textbf{Exceção:}
\begin{itemize}
    \item \textbf{Diagrama vazio:} Caso o diagrama não possua nenhuma tabela cadastrada, será exibida uma mensagem informando que não há elementos para exportar.
    \item \textbf{Erro de instabilidade dos servidores:} Caso algum erro ocorra nos servidores, será apresentada uma mensagem explicando que houve alguma falha de comunicação no sistema e encorajando o usuário a tentar novamente mais tarde.
\end{itemize}

\textbf{Pós-Condições:}
\begin{itemize}
    \item O Editor possui um arquivo de script DDL SQL correspondente ao diagrama ER do projeto.
\end{itemize}

\vspace{1em}

\begin{itemize}
    \item \textbf{CSU-24: Importar SQL}
\end{itemize}

\textbf{Caso de uso:} Importar SQL

\textbf{Identificador:} CSU-24

\textbf{Atores envolvidos:} Editor

\textbf{Sumário:} Esta funcionalidade permite que o Editor importe um script SQL/DDL para popular o diagrama ER do projeto de forma automática.

\textbf{Pré-condição:}
\begin{itemize}
    \item O usuário deverá estar em modo de edição do diagrama (CSU-07).
\end{itemize}

\textbf{Fluxo principal:}
\begin{itemize}
    \item O Editor seleciona a opção de importar um script SQL na interface de edição;
    \item O Editor fornece o script SQL/DDL, seja por meio de \textit{upload} de arquivo ou colagem direta do conteúdo;
    \item O sistema faz o \textit{parse} do script e gera as tabelas e relacionamentos correspondentes no diagrama;
    \item O diagrama é atualizado com a estrutura importada;
    \item As alterações são persistidas.
\end{itemize}

\textbf{Exceção:}
\begin{itemize}
    \item \textbf{Script SQL inválido:} Caso o script fornecido contenha erros de sintaxe ou estrutura não suportada, será exibida uma mensagem de erro indicando o problema encontrado.
    \item \textbf{Erro de instabilidade dos servidores:} Caso algum erro ocorra nos servidores, será apresentada uma mensagem explicando que houve alguma falha de comunicação no sistema e encorajando o usuário a tentar novamente mais tarde.
\end{itemize}

\textbf{Pós-Condições:}
\begin{itemize}
    \item O diagrama ER do projeto estará populado com as tabelas e relacionamentos definidos no script SQL importado.
\end{itemize}
